\section{\iflanguage{english}{Discussion}{Diskussion}}\label{sec:discussion}

This chapter discusses the results for each of the three questions of
\autoref{sec:motivation}, the limitations of the evidence, and open questions.

\subsection{Image--Label Fusion}\label{sec:discussion-fusion}

Replacing Medverse's two U-Nets with a shared encoder improved accuracy on the
training vocabulary by 0.16 to 0.23 Dice, but not on new targets: there,
single-axis attention was below Dual U-Net Cross-Attn on all eight datasets.
Dual attention improved over single-axis attention on 12 of 13 rows, most on
new targets. Letting image and label tokens interact in every layer thus seems
to help most when the model has to extract the target from the context masks
rather than recognize a trained class, although the results do not test this
mechanism directly.

Against Dual U-Net Cross-Attn, dual attention led on the held-out mean, but
this lead came from FLARE22 and TotalSeg MRI, whose classes are largely in the
training vocabulary. On the eight datasets with new targets, Dual U-Net
Cross-Attn was better on six. Its initialization from Medverse's weights,
trained on different data, may explain part of this; training both from the
same initialization would separate the two. Dual attention reached its results
with the fewest parameters of the three variants.

\subsection{Coarse-to-Fine Cascade}\label{sec:discussion-cascade}

Placing each crop from the previous level's prediction works: Dice increases
from level to level on every dataset (\autoref{tab:cascade-levels}), and on
TotalSeg CT the coarse prediction places the middle-level crop almost as well
as the ground truth does (92\% vs.\ 97\% of the target's voxels contained).

The clearest advantage over Medverse is speed: Cascade is more than an order
of magnitude faster. Since Medverse is faster at a single level, this comes
from processing one crop per level instead of the whole volume, not from the
architecture. The accuracy comparison is less clear. On the non-lesion datasets, Cascade leads on four of six against step-matched
Medverse but only on three of six at Medverse's native depth. Its largest
leads are on datasets from its training vocabulary, whereas Medverse leads on
MSD Hippocampus and NasalSeg, which neither model saw during training.

Training with the model's own prediction as prior was better than training
with a perturbed ground-truth mask, both on the training vocabulary and on the
three non-lesion datasets outside it. A likely reason is that the model sees
its own errors during training, which a hand-designed perturbation only
approximates. Medverse, trained with a ground-truth prior, shows the opposite
case: inside our cascade, passing on its own coarse prediction lowered its
Dice on six of seven datasets. Whether a prior helps at all compared with none
remains untested.

Region restriction makes the cost of each level independent of the volume
size, but makes the fine levels depend on the coarse localization. On
GNC\_705, the single-level Dice of 0.223 fell to 0.035 under cascade
inference, likely because the coarse level misses the small lesions. Small,
multi-focal, or low-contrast targets are most at risk. Larger margins, several
crops per volume, or a fallback to full-volume processing when the coarse
prediction is uncertain could address this.

\subsection{Synthetic Task Generation}\label{sec:discussion-synth}

With the synthetic-bank canvas, neither the mix of shape families nor
$p_\text{synth}$ affected the lesion datasets. Replacing the bank with real scans improved three of the four without reducing accuracy on
TotalSegmentator. What the shapes are painted on thus mattered more than which
shapes are painted. This change makes the canvas both real and, so the
two factors are not separated.

Under cascade inference, Synth + Cascade improved mainly on MSD Hippocampus
and NasalSeg, not on the lesion datasets that motivated the generator, and it
differs from Cascade in more than the synthetic tasks. Lesion Dice stayed low
throughout: at most 0.223 at a single level, and at most 0.070 for our models
under cascade inference. Lesions vary in location, number, and shape far more
than anatomical structures, so a single context example says little about
where the target is. In-context lesion segmentation remains largely unsolved
in this setting.

\subsection{Limitations}\label{sec:discussion-limitations}

\paragraph{Missing ablations.}
The cascade has not been evaluated without a prior or with full-volume
processing at the fine levels, so the contributions of the prior and of region
restriction are not isolated.

\paragraph{Evaluation.}
All results use a single context pair ($K = 1$), and the fusion comparison
uses crops centered on the ground truth. Medverse is the only baseline
evaluated empirically; Iris and UniverSeg are compared only by design.
TotalSeg MRI is part of our training data, and Medverse's training data may
overlap with MSD Hippocampus and NasalSeg, which affects the comparison on
these datasets.

\subsection{Open Questions}\label{sec:discussion-open}

First, the position axis attends over all $K+1$ volumes at once, so its cost
grows quadratically with the number of context volumes; how accuracy and cost
scale with $K$ is untested. Second, region restriction needs a way to handle
targets the coarse level cannot localize. Third, synthetic tasks did not yet
help on lesions under cascade inference; tasks that resemble lesions in
appearance and distribution, not only in shape, are one direction to explore.